% Options for packages loaded elsewhere
\PassOptionsToPackage{unicode}{hyperref}
\PassOptionsToPackage{hyphens}{url}
\PassOptionsToPackage{dvipsnames,svgnames*,x11names*}{xcolor}
%
\documentclass[
  9.5pt,
]{article}
\usepackage{amsmath,amssymb}
\usepackage{lmodern}
\usepackage{iftex}
\ifPDFTeX
  \usepackage[T1]{fontenc}
  \usepackage[utf8]{inputenc}
  \usepackage{textcomp} % provide euro and other symbols
\else % if luatex or xetex
  \usepackage{unicode-math}
  \defaultfontfeatures{Scale=MatchLowercase}
  \defaultfontfeatures[\rmfamily]{Ligatures=TeX,Scale=1}
  \setmainfont[]{Charter}
  \setmonofont[Scale=0.72]{Menlo}
\fi
% Use upquote if available, for straight quotes in verbatim environments
\IfFileExists{upquote.sty}{\usepackage{upquote}}{}
\IfFileExists{microtype.sty}{% use microtype if available
  \usepackage[]{microtype}
  \UseMicrotypeSet[protrusion]{basicmath} % disable protrusion for tt fonts
}{}
\makeatletter
\@ifundefined{KOMAClassName}{% if non-KOMA class
  \IfFileExists{parskip.sty}{%
    \usepackage{parskip}
  }{% else
    \setlength{\parindent}{0pt}
    \setlength{\parskip}{6pt plus 2pt minus 1pt}}
}{% if KOMA class
  \KOMAoptions{parskip=half}}
\makeatother
\usepackage{xcolor}
\IfFileExists{xurl.sty}{\usepackage{xurl}}{} % add URL line breaks if available
\IfFileExists{bookmark.sty}{\usepackage{bookmark}}{\usepackage{hyperref}}
\hypersetup{
  colorlinks=true,
  linkcolor={RoyalBlue},
  filecolor={Maroon},
  citecolor={Blue},
  urlcolor={RoyalBlue},
  pdfcreator={LaTeX via pandoc}}
\urlstyle{same} % disable monospaced font for URLs
\usepackage[margin=0.85in]{geometry}
\usepackage{color}
\usepackage{fancyvrb}
\newcommand{\VerbBar}{|}
\newcommand{\VERB}{\Verb[commandchars=\\\{\}]}
\DefineVerbatimEnvironment{Highlighting}{Verbatim}{commandchars=\\\{\}}
% Add ',fontsize=\small' for more characters per line
\usepackage{framed}
\definecolor{shadecolor}{RGB}{248,248,248}
\newenvironment{Shaded}{\begin{snugshade}}{\end{snugshade}}
\newcommand{\AlertTok}[1]{\textcolor[rgb]{0.94,0.16,0.16}{#1}}
\newcommand{\AnnotationTok}[1]{\textcolor[rgb]{0.56,0.35,0.01}{\textbf{\textit{#1}}}}
\newcommand{\AttributeTok}[1]{\textcolor[rgb]{0.77,0.63,0.00}{#1}}
\newcommand{\BaseNTok}[1]{\textcolor[rgb]{0.00,0.00,0.81}{#1}}
\newcommand{\BuiltInTok}[1]{#1}
\newcommand{\CharTok}[1]{\textcolor[rgb]{0.31,0.60,0.02}{#1}}
\newcommand{\CommentTok}[1]{\textcolor[rgb]{0.56,0.35,0.01}{\textit{#1}}}
\newcommand{\CommentVarTok}[1]{\textcolor[rgb]{0.56,0.35,0.01}{\textbf{\textit{#1}}}}
\newcommand{\ConstantTok}[1]{\textcolor[rgb]{0.00,0.00,0.00}{#1}}
\newcommand{\ControlFlowTok}[1]{\textcolor[rgb]{0.13,0.29,0.53}{\textbf{#1}}}
\newcommand{\DataTypeTok}[1]{\textcolor[rgb]{0.13,0.29,0.53}{#1}}
\newcommand{\DecValTok}[1]{\textcolor[rgb]{0.00,0.00,0.81}{#1}}
\newcommand{\DocumentationTok}[1]{\textcolor[rgb]{0.56,0.35,0.01}{\textbf{\textit{#1}}}}
\newcommand{\ErrorTok}[1]{\textcolor[rgb]{0.64,0.00,0.00}{\textbf{#1}}}
\newcommand{\ExtensionTok}[1]{#1}
\newcommand{\FloatTok}[1]{\textcolor[rgb]{0.00,0.00,0.81}{#1}}
\newcommand{\FunctionTok}[1]{\textcolor[rgb]{0.00,0.00,0.00}{#1}}
\newcommand{\ImportTok}[1]{#1}
\newcommand{\InformationTok}[1]{\textcolor[rgb]{0.56,0.35,0.01}{\textbf{\textit{#1}}}}
\newcommand{\KeywordTok}[1]{\textcolor[rgb]{0.13,0.29,0.53}{\textbf{#1}}}
\newcommand{\NormalTok}[1]{#1}
\newcommand{\OperatorTok}[1]{\textcolor[rgb]{0.81,0.36,0.00}{\textbf{#1}}}
\newcommand{\OtherTok}[1]{\textcolor[rgb]{0.56,0.35,0.01}{#1}}
\newcommand{\PreprocessorTok}[1]{\textcolor[rgb]{0.56,0.35,0.01}{\textit{#1}}}
\newcommand{\RegionMarkerTok}[1]{#1}
\newcommand{\SpecialCharTok}[1]{\textcolor[rgb]{0.00,0.00,0.00}{#1}}
\newcommand{\SpecialStringTok}[1]{\textcolor[rgb]{0.31,0.60,0.02}{#1}}
\newcommand{\StringTok}[1]{\textcolor[rgb]{0.31,0.60,0.02}{#1}}
\newcommand{\VariableTok}[1]{\textcolor[rgb]{0.00,0.00,0.00}{#1}}
\newcommand{\VerbatimStringTok}[1]{\textcolor[rgb]{0.31,0.60,0.02}{#1}}
\newcommand{\WarningTok}[1]{\textcolor[rgb]{0.56,0.35,0.01}{\textbf{\textit{#1}}}}
\usepackage{longtable,booktabs,array}
\usepackage{calc} % for calculating minipage widths
% Correct order of tables after \paragraph or \subparagraph
\usepackage{etoolbox}
\makeatletter
\patchcmd\longtable{\par}{\if@noskipsec\mbox{}\fi\par}{}{}
\makeatother
% Allow footnotes in longtable head/foot
\IfFileExists{footnotehyper.sty}{\usepackage{footnotehyper}}{\usepackage{footnote}}
\makesavenoteenv{longtable}
\setlength{\emergencystretch}{3em} % prevent overfull lines
\providecommand{\tightlist}{%
  \setlength{\itemsep}{0pt}\setlength{\parskip}{0pt}}
\setcounter{secnumdepth}{-\maxdimen} % remove section numbering
\usepackage{newunicodechar}
\newunicodechar{→}{\ensuremath{\rightarrow}}
\newunicodechar{←}{\ensuremath{\leftarrow}}
\newunicodechar{↔}{\ensuremath{\leftrightarrow}}
\newunicodechar{⇒}{\ensuremath{\Rightarrow}}
\newunicodechar{≈}{\ensuremath{\approx}}
\newunicodechar{≤}{\ensuremath{\leq}}
\newunicodechar{≥}{\ensuremath{\geq}}
\newunicodechar{≠}{\ensuremath{\neq}}
\newunicodechar{×}{\ensuremath{\times}}
\newunicodechar{±}{\ensuremath{\pm}}
\newunicodechar{∈}{\ensuremath{\in}}
\newunicodechar{≪}{\ensuremath{\ll}}
\ifLuaTeX
  \usepackage{selnolig}  % disable illegal ligatures
\fi

\title{SDR Platform — Experiment Operating Guide}
\author{}
\date{}

\begin{document}
\maketitle

{
\hypersetup{linkcolor=}
\setcounter{tocdepth}{2}
\tableofcontents
}
\hypertarget{sdr-platform--experiment-operating-guide}{%
\section{SDR Platform --- Experiment Operating
Guide}\label{sdr-platform--experiment-operating-guide}}

A step-by-step runbook for operating \textbf{every application} on the
shared SDR PHY. Each application has a \textbf{radio-free} path
(validate the logic with no hardware) and a \textbf{hardware} path (two
hosts + USRPs). Do the radio-free steps first --- they catch 90\% of
mistakes before a radio is involved.

\begin{longtable}[]{@{}lllll@{}}
\toprule
App & Folder & Built? & Radio-free & Hardware \\
\midrule
\endhead
1. Reliable link --- \textbf{MARL / OSU} & \texttt{MARL\_RA\_Union/},
\texttt{python/} & yes & yes & yes (needs ≥3 USRP for real
collisions) \\
1. Reliable link --- \textbf{Federated Learning} & \texttt{python/fl.py}
& yes & yes & yes \\
2. \textbf{STC-AirComp} & \texttt{STC\_AirComp\_Union/} & design &
planned & planned \\
3. \textbf{CLIP semantic comm} & \texttt{CLIP\_SemCom\_Union/} & yes &
yes & yes \\
\bottomrule
\end{longtable}

Conventions in this guide: \texttt{\$AP\_IP} / \texttt{\$SERVER\_IP} /
\texttt{\$RX\_IP} = the receiver/AP host's LAN IP;
\texttt{serial=30CD424} and \texttt{serial=30CD3F7} = the two B210s;
\texttt{addr=192.168.20.2} = the N210. \textbf{Always start the
receiver/AP first}, then the transmitter.

\begin{center}\rule{0.5\linewidth}{0.5pt}\end{center}

\hypertarget{0-one-time-setup-do-once-per-host}{%
\subsection{0. One-time setup (do once per
host)}\label{0-one-time-setup-do-once-per-host}}

\hypertarget{01-install-dependencies}{%
\subsubsection{0.1 Install dependencies}\label{01-install-dependencies}}

\begin{Shaded}
\begin{Highlighting}[]
\BuiltInTok{cd}\NormalTok{ Hardware\_update}
\ExtensionTok{./deploy/initialization.sh}            \CommentTok{\# C++ PHY deps (UHD, Boost, FFTW3f, VOLK) + Python (numpy, matplotlib)}
\ExtensionTok{./deploy/initialization.sh} \AttributeTok{{-}{-}build}    \CommentTok{\# also configure + compile build/sdr\_system}
\end{Highlighting}
\end{Shaded}

\hypertarget{02-build-the-c-phy-sdr_system--natively-on-each-box}{%
\subsubsection{\texorpdfstring{0.2 Build the C++ PHY
(\texttt{sdr\_system}) --- natively on each
box}{0.2 Build the C++ PHY (sdr\_system) --- natively on each box}}\label{02-build-the-c-phy-sdr_system--natively-on-each-box}}

The binary is \textbf{architecture-specific} --- a Mac (arm64/x86\_64)
binary will not run on the Linux lab host
(\texttt{Exec\ format\ error}). Build on the machine that will run it:

\begin{Shaded}
\begin{Highlighting}[]
\BuiltInTok{cd}\NormalTok{ Hardware\_update/phy}
\ExtensionTok{conda}\NormalTok{ deactivate }\DecValTok{2}\OperatorTok{\textgreater{}}\NormalTok{/dev/null }\KeywordTok{||} \FunctionTok{true}      \CommentTok{\# Anaconda\textquotesingle{}s cmake bakes a dead path and breaks \textasciigrave{}make\textasciigrave{}}
\FunctionTok{rm} \AttributeTok{{-}rf}\NormalTok{ build }\KeywordTok{\&\&} \FunctionTok{mkdir}\NormalTok{ build }\KeywordTok{\&\&} \BuiltInTok{cd}\NormalTok{ build}
\FunctionTok{cmake}\NormalTok{ ..            }\CommentTok{\# if \textasciigrave{}cmake\textasciigrave{} is Anaconda\textquotesingle{}s, use /usr/bin/cmake .. explicitly}
\FunctionTok{make} \AttributeTok{{-}j}
\FunctionTok{ls}\NormalTok{ ./sdr\_system     }\CommentTok{\# \textless{}{-} the binary}
\end{Highlighting}
\end{Shaded}

Rebuilding also regenerates \texttt{phy/python/sdr.py} and
\texttt{PARAMETERS.md} (CMake post-build).

\hypertarget{03-build-the-block-api-pyphy-needed-for-the-clip-pyphy-channel-optional-otherwise}{%
\subsubsection{\texorpdfstring{0.3 Build the block API \texttt{pyphy}
(needed for the CLIP \texttt{pyphy} channel; optional
otherwise)}{0.3 Build the block API pyphy (needed for the CLIP pyphy channel; optional otherwise)}}\label{03-build-the-block-api-pyphy-needed-for-the-clip-pyphy-channel-optional-otherwise}}

\begin{Shaded}
\begin{Highlighting}[]
\BuiltInTok{cd}\NormalTok{ Hardware\_update/phy}
\ExtensionTok{bindings/build.sh}                 \CommentTok{\# DSP blocks only (builds anywhere)}
\VariableTok{WITH\_UHD}\OperatorTok{=}\NormalTok{1 }\ExtensionTok{bindings/build.sh}      \CommentTok{\# + the Radio source/sink (on the lab host, needs UHD)}
\CommentTok{\# On macOS run pyphy code as:  PYTHONPATH=phy/bindings arch {-}x86\_64 python3 ...}
\end{Highlighting}
\end{Shaded}

\hypertarget{04-confirm-the-radios-are-visible}{%
\subsubsection{0.4 Confirm the radios are
visible}\label{04-confirm-the-radios-are-visible}}

\begin{Shaded}
\begin{Highlighting}[]
\ExtensionTok{uhd\_find\_devices}                              \CommentTok{\# lists connected USRPs}
\ExtensionTok{uhd\_find\_devices} \AttributeTok{{-}{-}args}\NormalTok{ addr=192.168.20.2     }\CommentTok{\# the N210 specifically}
\end{Highlighting}
\end{Shaded}

\hypertarget{05-pick-a-clean-carrier-recommended-before-any-rf-run}{%
\subsubsection{0.5 Pick a clean carrier (recommended before any RF
run)}\label{05-pick-a-clean-carrier-recommended-before-any-rf-run}}

\begin{Shaded}
\begin{Highlighting}[]
\BuiltInTok{cd}\NormalTok{ phy/python}
\ExtensionTok{python3}\NormalTok{ freq\_scan.py }\AttributeTok{{-}{-}start}\NormalTok{ 902 }\AttributeTok{{-}{-}stop}\NormalTok{ 928 }\AttributeTok{{-}{-}step}\NormalTok{ 1 }\AttributeTok{{-}{-}rx{-}args}\NormalTok{ addr=192.168.20.2}
\CommentTok{\# use the quietest highlighted frequency as your link freq (default is 915 MHz)}
\end{Highlighting}
\end{Shaded}

\hypertarget{06-golden-rules-save-yourself-hours}{%
\subsubsection{0.6 Golden rules (save yourself
hours)}\label{06-golden-rules-save-yourself-hours}}

\begin{itemize}
\tightlist
\item
  \textbf{RX/AP first, TX second.} Keep both radios \textbf{warm}
  (started once, left running).
\item
  \textbf{N210 rates are exact:}
  \texttt{-\/-rx-rate\ 2e6\ -\/-symbol-rate\ 1e6} (1.6e6 does not snap
  on its 100 MHz clock).
\item
  \textbf{Two-host ACK:} each side's \texttt{-\/-ack-host} = the
  \emph{peer} host's IP.
\item
  \textbf{Free-running LOs:} single-carrier coherent QPSK ≈ 17\%
  delivery. Use \textbf{DQPSK} (single-carrier) or \textbf{OFDM}
  (coherent QPSK, pilots track CFO). A shared 10 MHz reference makes
  coherent reliable.
\item
  \textbf{Stuck RX won't Ctrl-C?} Press Ctrl-C again (escalating
  handler); emergency kill = \texttt{Ctrl-\textbackslash{}} or
  \texttt{kill\ -9}.
\item
  \textbf{\texttt{\textless{}frozen\ getpath\textgreater{}} crash} =
  stale shell CWD -\textgreater{} \texttt{cd} out and back in.
\end{itemize}

\begin{center}\rule{0.5\linewidth}{0.5pt}\end{center}

\hypertarget{1a-application-1--marl--osu-random-access}{%
\subsection{1A. Application 1 --- MARL / OSU random
access}\label{1a-application-1--marl--osu-random-access}}

RL agents share one channel; the \textbf{ACK is the reward} (delivered =
success, timeout = collision/loss). Single-shot bursts
(\texttt{-\/-max-attempts\ 1}); the \emph{policy} owns retransmission.

\hypertarget{step-0--radio-free-validation-no-hardware}{%
\subsubsection{Step 0 --- radio-free validation (no
hardware)}\label{step-0--radio-free-validation-no-hardware}}

\begin{Shaded}
\begin{Highlighting}[]
\BuiltInTok{cd}\NormalTok{ applications/MARL\_RA\_Union}
\ExtensionTok{python3}\NormalTok{ ap\_multi.py }\AttributeTok{{-}{-}self{-}test}      \CommentTok{\# per{-}agent ACK routing (synthetic id stream)  {-}\textgreater{} PASS}
\ExtensionTok{python3}\NormalTok{ ap\_multi.py }\AttributeTok{{-}{-}sim{-}test}       \CommentTok{\# parser + end{-}to{-}end with a fake C++ sink       {-}\textgreater{} PASS}
\ExtensionTok{python3}\NormalTok{ slot\_sync.py }\AttributeTok{{-}{-}self{-}test}     \CommentTok{\# slot clock aligns 3 clients                     {-}\textgreater{} PASS}

\CommentTok{\# offline decentralized run (one process per agent, mock medium):}
\ExtensionTok{python3}\NormalTok{ mock\_medium.py }\AttributeTok{{-}{-}agents}\NormalTok{ 2 }\AttributeTok{{-}{-}slots}\NormalTok{ 600 }\KeywordTok{\&}
\ExtensionTok{python3}\NormalTok{ agent\_node.py }\AttributeTok{{-}{-}mock} \AttributeTok{{-}{-}id}\NormalTok{ 0 }\AttributeTok{{-}{-}slots}\NormalTok{ 600}
\ExtensionTok{python3}\NormalTok{ agent\_node.py }\AttributeTok{{-}{-}mock} \AttributeTok{{-}{-}id}\NormalTok{ 1 }\AttributeTok{{-}{-}slots}\NormalTok{ 600}

\CommentTok{\# offline training (learns to transmit; prints P(transmit|queued) climbing):}
\ExtensionTok{python3}\NormalTok{ marl\_train.py }\AttributeTok{{-}{-}mock} \AttributeTok{{-}{-}steps}\NormalTok{ 400                       }\CommentTok{\# single agent (A2C)}
\ExtensionTok{python3}\NormalTok{ marl\_multi\_train.py }\AttributeTok{{-}{-}mock} \AttributeTok{{-}{-}agents}\NormalTok{ 4 }\AttributeTok{{-}{-}steps}\NormalTok{ 800 }\DataTypeTok{\textbackslash{}}
        \AttributeTok{{-}{-}coll{-}penalty}\NormalTok{ 0.5 }\AttributeTok{{-}{-}compare{-}aloha} \StringTok{"0.15,0.25,0.5"}     \CommentTok{\# multi{-}agent vs ALOHA}
\end{Highlighting}
\end{Shaded}

\textbf{Success:} self/sim-tests PASS; in training
\texttt{P(transmit\textbar{}queued)} rises on a clean mock link;
multi-agent collision rate drops and per-agent \texttt{P(transmit)}
settles near \texttt{1/N}.

\hypertarget{step-1--hardware-single-agent-2-usrps-1-ap--1-agent}{%
\subsubsection{Step 1 --- hardware, single agent (2 USRPs: 1 AP + 1
agent)}\label{step-1--hardware-single-agent-2-usrps-1-ap--1-agent}}

\begin{Shaded}
\begin{Highlighting}[]
\CommentTok{\# AP host (warm receiver + ACK router):}
\ExtensionTok{python3}\NormalTok{ real\_channel.py ap }\AttributeTok{{-}{-}rx{-}args}\NormalTok{ serial=30CD3F7}

\CommentTok{\# Agent host (trains online from real ACK/timeout):}
\ExtensionTok{python3}\NormalTok{ marl\_train.py }\AttributeTok{{-}{-}tx{-}args}\NormalTok{ serial=30CD424 }\AttributeTok{{-}{-}steps}\NormalTok{ 150 }\AttributeTok{{-}{-}scheme}\NormalTok{ DQPSK}
\end{Highlighting}
\end{Shaded}

\textbf{Success:} the agent fires real bursts; reward tracks real
delivery; the model + reward plot save next to \texttt{-\/-out}. (One
agent has no collision partner --- use Step 2 or the jammer for
contention.)

\hypertarget{step-2--hardware-decentralized-multi-agent-3-usrps-1-ap--2-agents}{%
\subsubsection{Step 2 --- hardware, decentralized multi-agent (3 USRPs:
1 AP + 2
agents)}\label{step-2--hardware-decentralized-multi-agent-3-usrps-1-ap--2-agents}}

\begin{Shaded}
\begin{Highlighting}[]
\CommentTok{\# AP host — start BOTH the decoder/ACK router and the slot clock:}
\ExtensionTok{python3}\NormalTok{ ap\_multi.py }\AttributeTok{{-}{-}agents}\NormalTok{ 2 }\AttributeTok{{-}{-}scheme}\NormalTok{ QPSK }\AttributeTok{{-}{-}rx{-}args}\NormalTok{ serial=30CD3F7    }\CommentTok{\# B210 AP}
\CommentTok{\#   (N210 AP instead:  {-}{-}rx{-}args addr=192.168.20.2 {-}{-}rx{-}subdev A:0 {-}{-}rx{-}rate 2e6 {-}{-}symbol{-}rate 1e6 {-}{-}scheme DQPSK)}
\ExtensionTok{python3}\NormalTok{ slot\_sync.py }\AttributeTok{{-}{-}agents}\NormalTok{ 2 }\AttributeTok{{-}{-}slot{-}ms}\NormalTok{ 150                            }\CommentTok{\# slot clock (port 5600)}

\CommentTok{\# Each agent host (local{-}only obs; learned CSMA), pointing at the AP host:}
\ExtensionTok{python3}\NormalTok{ agent\_node.py }\AttributeTok{{-}{-}id}\NormalTok{ 0 }\AttributeTok{{-}{-}tx{-}args}\NormalTok{ serial=30CD424 }\AttributeTok{{-}{-}scheme}\NormalTok{ QPSK }\AttributeTok{{-}{-}ap{-}host} \VariableTok{$AP\_IP} \AttributeTok{{-}{-}slot{-}host} \VariableTok{$AP\_IP}
\ExtensionTok{python3}\NormalTok{ agent\_node.py }\AttributeTok{{-}{-}id}\NormalTok{ 1 }\AttributeTok{{-}{-}tx{-}args}\NormalTok{ serial=}\OperatorTok{\textless{}}\NormalTok{3RD}\OperatorTok{\textgreater{}}\NormalTok{   {-}{-}scheme QPSK }\AttributeTok{{-}{-}ap{-}host} \VariableTok{$AP\_IP} \AttributeTok{{-}{-}slot{-}host} \VariableTok{$AP\_IP}
\end{Highlighting}
\end{Shaded}

\textbf{Success:} AP prints
\texttt{{[}BURST{]}\ id=\ldots{}\ hex=\ldots{}} per decoded frame and
routes an ACK to that agent; ≥2 intents in a slot ⇒ COLLISION (no ACK);
agents learn to back off (collision rate falls).

\hypertarget{optional--add-real-contention-with-the-jammer}{%
\subsubsection{Optional --- add real contention with the
jammer}\label{optional--add-real-contention-with-the-jammer}}

\begin{Shaded}
\begin{Highlighting}[]
\ExtensionTok{python3}\NormalTok{ ../jammer/jammer.py }\AttributeTok{{-}{-}tx{-}args}\NormalTok{ serial=}\OperatorTok{\textless{}}\NormalTok{JAMMER}\OperatorTok{\textgreater{}}\NormalTok{ {-}{-}freq 915e6 }\DataTypeTok{\textbackslash{}}
        \AttributeTok{{-}{-}mode}\NormalTok{ burst }\AttributeTok{{-}{-}interval}\NormalTok{ 150 }\AttributeTok{{-}{-}duration}\NormalTok{ 300 }\AttributeTok{{-}{-}tx{-}gain}\NormalTok{ 80}
\end{Highlighting}
\end{Shaded}

\begin{center}\rule{0.5\linewidth}{0.5pt}\end{center}

\hypertarget{1b-application-1--federated-learning-mnist}{%
\subsection{1B. Application 1 --- Federated Learning
(MNIST)}\label{1b-application-1--federated-learning-mnist}}

Clients train locally and send \textbf{compressed model deltas} (float32
vectors) up to a server that FedAvg-aggregates and broadcasts the
aggregate down.

\hypertarget{step-0--radio-free}{%
\subsubsection{Step 0 --- radio-free}\label{step-0--radio-free}}

\begin{Shaded}
\begin{Highlighting}[]
\BuiltInTok{cd}\NormalTok{ applications/FL\_Union}
\ExtensionTok{python3}\NormalTok{ fl.py }\AttributeTok{{-}{-}mock} \AttributeTok{{-}{-}clients}\NormalTok{ 2 }\AttributeTok{{-}{-}rounds}\NormalTok{ 20                    }\CommentTok{\# compressed (default) {-}\textgreater{} \textasciitilde{}0.71{-}\textgreater{}0.93}
\ExtensionTok{python3}\NormalTok{ fl.py }\AttributeTok{{-}{-}mock} \AttributeTok{{-}{-}clients}\NormalTok{ 2 }\AttributeTok{{-}{-}rounds}\NormalTok{ 20 }\AttributeTok{{-}{-}compress{-}ratio}\NormalTok{ 0 }\CommentTok{\# full model}

\CommentTok{\# whole protocol over TCP (no radios), real server/client processes:}
\ExtensionTok{python3}\NormalTok{ fl.py server }\AttributeTok{{-}{-}clients}\NormalTok{ 1 }\AttributeTok{{-}{-}rounds}\NormalTok{ 20 }\AttributeTok{{-}{-}uplink}\NormalTok{ tcp }\AttributeTok{{-}{-}downlink}\NormalTok{ tcp }\KeywordTok{\&}
\ExtensionTok{python3}\NormalTok{ fl.py client }\AttributeTok{{-}{-}client{-}id}\NormalTok{ 0 }\AttributeTok{{-}{-}clients}\NormalTok{ 1 }\AttributeTok{{-}{-}rounds}\NormalTok{ 20 }\AttributeTok{{-}{-}uplink}\NormalTok{ tcp }\AttributeTok{{-}{-}downlink}\NormalTok{ tcp }\DataTypeTok{\textbackslash{}}
        \AttributeTok{{-}{-}net{-}host}\NormalTok{ 127.0.0.1 }\AttributeTok{{-}{-}net{-}port}\NormalTok{ 5700 }\AttributeTok{{-}{-}ack{-}host}\NormalTok{ 127.0.0.1}
\end{Highlighting}
\end{Shaded}

\textbf{Success:} mock accuracy climbs
\textasciitilde0.71-\textgreater0.93 in 20 rounds; TCP run reaches
\textasciitilde0.92 in 12 rounds.

\hypertarget{step-1--hardware-server--n210-rx-only-client--b210-tx-start-the-server-first}{%
\subsubsection{Step 1 --- hardware (server = N210 RX-only, client = B210
TX). Start the server
first.}\label{step-1--hardware-server--n210-rx-only-client--b210-tx-start-the-server-first}}

\begin{Shaded}
\begin{Highlighting}[]
\CommentTok{\# Server host (N210): uplink over RF, model broadcast back over TCP}
\ExtensionTok{python3}\NormalTok{ fl.py server }\AttributeTok{{-}{-}clients}\NormalTok{ 1 }\AttributeTok{{-}{-}rounds}\NormalTok{ 20 }\AttributeTok{{-}{-}uplink}\NormalTok{ wireless }\AttributeTok{{-}{-}downlink}\NormalTok{ tcp }\DataTypeTok{\textbackslash{}}
        \AttributeTok{{-}{-}rx{-}args}\NormalTok{ addr=192.168.20.2 }\AttributeTok{{-}{-}rx{-}subdev}\NormalTok{ A:0 }\AttributeTok{{-}{-}net{-}port}\NormalTok{ 5700}

\CommentTok{\# Client host (B210):  ack{-}host + net{-}host = the server\textquotesingle{}s IP}
\ExtensionTok{python3}\NormalTok{ fl.py client }\AttributeTok{{-}{-}client{-}id}\NormalTok{ 0 }\AttributeTok{{-}{-}clients}\NormalTok{ 1 }\AttributeTok{{-}{-}rounds}\NormalTok{ 20 }\AttributeTok{{-}{-}uplink}\NormalTok{ wireless }\AttributeTok{{-}{-}downlink}\NormalTok{ tcp }\DataTypeTok{\textbackslash{}}
        \AttributeTok{{-}{-}tx{-}args}\NormalTok{ serial=30CD424 }\AttributeTok{{-}{-}ack{-}host} \VariableTok{$SERVER\_IP} \AttributeTok{{-}{-}net{-}host} \VariableTok{$SERVER\_IP} \AttributeTok{{-}{-}net{-}port}\NormalTok{ 5700}
\end{Highlighting}
\end{Shaded}

\textbf{Success:} each round the client's delta is delivered (ARQ), the
server aggregates and broadcasts; accuracy rises round over round. Use
\texttt{-\/-waveform\ ofdm} for coherent QPSK if the SC link is
marginal.

\begin{center}\rule{0.5\linewidth}{0.5pt}\end{center}

\hypertarget{2-application-2--stc-aircomp-design-only--not-yet-runnable}{%
\subsection{2. Application 2 --- STC-AirComp (design only --- not yet
runnable)}\label{2-application-2--stc-aircomp-design-only--not-yet-runnable}}

This app is specified but \textbf{not implemented} (blocked on the
paper's exact STLC precoding/combiner equations and on a synchronized
2-channel RX capture). See \texttt{STC\_AirComp\_Union/INTEGRATION.md}.
When built, the operating sequence will be:

\begin{enumerate}
\def\labelenumi{\arabic{enumi}.}
\tightlist
\item
  \textbf{DSP loopback (radio-free):} validate a 2-sensor
  -\textgreater{} 1-AP sum through a simulated 2-antenna AWGN channel (0
  error at high SNR; NMSE-vs-SNR curve).
\item
  \textbf{\texttt{capture2} bring-up:} verify a synchronized 2-channel
  RX on a B210/X310 (2×2) captures one known TX.
\item
  \textbf{Single-sensor RF:} one B210 sensor -\textgreater{} 2×2 AP;
  STLC-encode, transmit, \texttt{capture2}, combine, recover the single
  value. Measure the timing/phase coherence budget.
\item
  \textbf{Two-sensor RF sum:} add a second B210, \texttt{slot\_sync}
  them into one slot, recover \texttt{v1+v2}; report NMSE. Try
  \textbf{DBPSK} first (no shared clock), then BPSK +
  \texttt{-\/-ref\ external} if coherence demands it.
\end{enumerate}

To start building it, provide the paper's STLC matrices + estimator (see
the INTEGRATION note's §7).

\begin{center}\rule{0.5\linewidth}{0.5pt}\end{center}

\hypertarget{3-application-3--clip-semantic-communication}{%
\subsection{3. Application 3 --- CLIP semantic
communication}\label{3-application-3--clip-semantic-communication}}

The base station encodes an image with \textbf{CLIP} into a float32
embedding and transmits \textbf{only the embedding}; the user classifies
it (zero-shot) with no return link. Data-transfer archetype.

\begin{Shaded}
\begin{Highlighting}[]
\BuiltInTok{cd}\NormalTok{ applications/CLIP\_SemCom\_Union}
\end{Highlighting}
\end{Shaded}

\hypertarget{step-0--radio-free-no-torch-weights-needed-uses-the-mock-clip}{%
\subsubsection{Step 0 --- radio-free (no torch weights needed; uses the
mock
CLIP)}\label{step-0--radio-free-no-torch-weights-needed-uses-the-mock-clip}}

\begin{Shaded}
\begin{Highlighting}[]
\ExtensionTok{python3}\NormalTok{ semcom.py demo }\AttributeTok{{-}{-}mock}                       \CommentTok{\# encode {-}\textgreater{} lossless channel {-}\textgreater{} classify (\textasciitilde{}0.99)}

\CommentTok{\# reproduce the paper\textquotesingle{}s accuracy{-}vs{-}noise (Fig. 3) through OUR real modem + FEC:}
\VariableTok{PYTHONPATH}\OperatorTok{=}\NormalTok{../../phy/bindings }\FunctionTok{arch} \AttributeTok{{-}x86\_64}\NormalTok{ python3 semcom.py demo }\AttributeTok{{-}{-}mock} \DataTypeTok{\textbackslash{}}
        \AttributeTok{{-}{-}channel}\NormalTok{ pyphy }\AttributeTok{{-}{-}scheme}\NormalTok{ QPSK }\AttributeTok{{-}{-}fec}\NormalTok{ turbo }\AttributeTok{{-}{-}snr{-}sweep}\NormalTok{ 0,2,4,6,10}

\ExtensionTok{python3}\NormalTok{ semcom.py demo }\AttributeTok{{-}{-}mock} \AttributeTok{{-}{-}model{-}sweep}         \CommentTok{\# 3{-}CLIP{-}model accuracy / payload / delay / energy}
\end{Highlighting}
\end{Shaded}

\textbf{Success:} clean accuracy ≈ 1.0; in the \texttt{pyphy} sweep
accuracy climbs with SNR tracking BER (≈0 @0 dB -\textgreater{}
\textasciitilde0.7 @4 dB -\textgreater{} \textasciitilde0.99 @6 dB);
\texttt{int8} quant (\texttt{-\/-quant\ int8}) is 4× smaller payload.

\hypertarget{step-1--real-clip-weights-optional-torch-already-installed}{%
\subsubsection{Step 1 --- real CLIP weights (optional; torch already
installed)}\label{step-1--real-clip-weights-optional-torch-already-installed}}

\begin{Shaded}
\begin{Highlighting}[]
\ExtensionTok{pip}\NormalTok{ install }\AttributeTok{{-}r}\NormalTok{ requirements.txt        }\CommentTok{\# open\_clip\_torch, pillow, ftfy, regex}
\ExtensionTok{python3}\NormalTok{ semcom.py demo }\AttributeTok{{-}{-}model}\NormalTok{ vit{-}l14 }\CommentTok{\# drops {-}{-}mock automatically {-}\textgreater{} real ViT{-}L/14}
\end{Highlighting}
\end{Shaded}

\hypertarget{step-2--over-the-radio-rx-host--n210-tx-host--b210-start-rx-first}{%
\subsubsection{Step 2 --- over the radio (RX host = N210, TX host =
B210). Start RX
first.}\label{step-2--over-the-radio-rx-host--n210-tx-host--b210-start-rx-first}}

\begin{Shaded}
\begin{Highlighting}[]
\CommentTok{\# RX host (N210): receive embeddings and classify locally}
\ExtensionTok{python3}\NormalTok{ semcom.py rx }\AttributeTok{{-}{-}rx{-}args}\NormalTok{ addr=192.168.20.2 }\AttributeTok{{-}{-}rx{-}subdev}\NormalTok{ A:0 }\AttributeTok{{-}{-}num}\NormalTok{ 8}

\CommentTok{\# TX host (B210): encode images, send each embedding (ack{-}host = RX host IP)}
\ExtensionTok{python3}\NormalTok{ semcom.py tx }\AttributeTok{{-}{-}tx{-}args}\NormalTok{ serial=30CD424 }\AttributeTok{{-}{-}ack{-}host} \VariableTok{$RX\_IP}
\end{Highlighting}
\end{Shaded}

\textbf{Success:} the RX prints
\texttt{predicted=\textless{}class\textgreater{}\ (truth=\textless{}class\textgreater{})}
per embedding. The radio path uses reliable ARQ (lossless delivery). For
the on-air noise-\textgreater accuracy curve, use the
\texttt{-\/-channel\ pyphy} sweep in Step 0 (radio-free), or the phase-2
single-shot mode in the INTEGRATION note.

\begin{center}\rule{0.5\linewidth}{0.5pt}\end{center}

\hypertarget{4-troubleshooting-the-gotchas-that-actually-bite}{%
\subsection{4. Troubleshooting (the gotchas that actually
bite)}\label{4-troubleshooting-the-gotchas-that-actually-bite}}

\begin{longtable}[]{@{}ll@{}}
\toprule
Symptom & Cause -\textgreater{} Fix \\
\midrule
\endhead
\texttt{Exec\ format\ error} running \texttt{sdr\_system} & Wrong-arch
binary -\textgreater{} rebuild natively on that box (§0.2). \\
\texttt{\textless{}frozen\ getpath\textgreater{}} Python crash & Stale
shell CWD -\textgreater{} \texttt{cd} out and back in. \\
Sink fails the \textbf{rate check} at startup & The receive-only sink
still checks the full chain -\textgreater{} set
\texttt{tx\_rate\ =\ rx\_rate} on it; N210 needs
\texttt{-\/-rx-rate\ 2e6\ -\/-symbol-rate\ 1e6}. \\
\texttt{ACQ\ ERROR\ /\ Filtered\ vector\ too\ short} (works at small
\texttt{k}, fails at large) & Long frame's energy dips below the
detector -\textgreater{} raise TX power/gain, add \texttt{-\/-det-mult},
\texttt{-\/-energy\_packet\_size}, or use smaller
\texttt{-\/-ldpc-k}. \\
Single-shot burst gets \textbf{no ACK} (MARL) & Cold LO / free-running
CFO -\textgreater{} keep the AP \textbf{warm}, use \textbf{DQPSK},
ideally a shared 10 MHz clock; "no ACK" is then ≈ collision. \\
Coherent \textbf{QPSK \textasciitilde17\% delivery} & Free-running LO
-\textgreater{} use \textbf{DQPSK} (SC) or \textbf{OFDM} (coherent,
pilots track CFO). \\
CLIP accuracy \textasciitilde0 even at high SNR (uncoded) & Any bit
error wrecks the 512-float payload -\textgreater{} add FEC
(\texttt{-\/-fec\ turbo}), use soft decode. \\
RX won't stop on Ctrl-C & Press Ctrl-C again (escalating handler);
emergency = \texttt{Ctrl-\textbackslash{}} or \texttt{kill\ -9}. \\
Two processes on one host clash & Give distinct ports: ARQ
\texttt{-\/-ack-port}, FL downlink \texttt{-\/-net-port}, slot clock
\texttt{-\/-port\ 5600}. \\
pyphy \texttt{Python.h\ not\ found} / arch mismatch (macOS) & Use
\texttt{phy/bindings/build.sh} (adds \texttt{-isystem},
\texttt{python3-config}); run with \texttt{arch\ -x86\_64\ python3}. \\
\bottomrule
\end{longtable}

\begin{center}\rule{0.5\linewidth}{0.5pt}\end{center}

\hypertarget{5-quick-reference-card}{%
\subsection{5. Quick-reference card}\label{5-quick-reference-card}}

\begin{Shaded}
\begin{Highlighting}[]
\CommentTok{\# {-}{-}{-}{-} radio{-}free (no hardware) {-}{-}{-}{-}}
\KeywordTok{(}\BuiltInTok{cd}\NormalTok{ applications/MARL\_RA\_Union }\KeywordTok{\&\&} \ExtensionTok{python3}\NormalTok{ ap\_multi.py }\AttributeTok{{-}{-}sim{-}test}\KeywordTok{)}      \CommentTok{\# MARL routing}
\KeywordTok{(}\BuiltInTok{cd}\NormalTok{ applications/MARL\_RA\_Union }\KeywordTok{\&\&} \ExtensionTok{python3}\NormalTok{ marl\_multi\_train.py }\AttributeTok{{-}{-}mock} \AttributeTok{{-}{-}agents}\NormalTok{ 4 }\AttributeTok{{-}{-}steps}\NormalTok{ 800 }\AttributeTok{{-}{-}coll{-}penalty}\NormalTok{ 0.5}\KeywordTok{)}
\KeywordTok{(}\BuiltInTok{cd}\NormalTok{ applications/FL\_Union }\KeywordTok{\&\&} \ExtensionTok{python3}\NormalTok{ fl.py }\AttributeTok{{-}{-}mock} \AttributeTok{{-}{-}clients}\NormalTok{ 2 }\AttributeTok{{-}{-}rounds}\NormalTok{ 20}\KeywordTok{)}  \CommentTok{\# FL}
\BuiltInTok{cd}\NormalTok{ applications/CLIP\_SemCom\_Union }\KeywordTok{\&\&} \ExtensionTok{python3}\NormalTok{ semcom.py demo }\AttributeTok{{-}{-}mock}   \CommentTok{\# CLIP}
\VariableTok{PYTHONPATH}\OperatorTok{=}\NormalTok{../../phy/bindings }\FunctionTok{arch} \AttributeTok{{-}x86\_64}\NormalTok{ python3 semcom.py demo }\AttributeTok{{-}{-}mock} \AttributeTok{{-}{-}channel}\NormalTok{ pyphy }\AttributeTok{{-}{-}fec}\NormalTok{ turbo }\AttributeTok{{-}{-}snr{-}sweep}\NormalTok{ 0,2,4,6,10}

\CommentTok{\# {-}{-}{-}{-} hardware (RX/AP host first, TX host second) {-}{-}{-}{-}}
\CommentTok{\# MARL single:   real\_channel.py ap {-}{-}rx{-}args serial=30CD3F7   |   marl\_train.py {-}{-}tx{-}args serial=30CD424 {-}{-}steps 150}
\CommentTok{\# MARL multi:    ap\_multi.py {-}{-}agents 2 {-}{-}rx{-}args serial=30CD3F7 + slot\_sync.py {-}{-}agents 2 | agent\_node.py {-}{-}id N {-}{-}ap{-}host $AP\_IP {-}{-}slot{-}host $AP\_IP}
\CommentTok{\# FL:            fl.py server {-}{-}uplink wireless {-}{-}downlink tcp {-}{-}rx{-}args addr=192.168.20.2 {-}{-}rx{-}subdev A:0 | fl.py client {-}{-}tx{-}args serial=30CD424 {-}{-}ack{-}host $SERVER\_IP {-}{-}net{-}host $SERVER\_IP}
\CommentTok{\# CLIP:          semcom.py rx {-}{-}rx{-}args addr=192.168.20.2 {-}{-}rx{-}subdev A:0 | semcom.py tx {-}{-}tx{-}args serial=30CD424 {-}{-}ack{-}host $RX\_IP}
\end{Highlighting}
\end{Shaded}

For engine options: \texttt{PARAMETERS.md} (auto-generated, complete).
For the full PHY reference: \texttt{../SYSTEM\_REFERENCE.pdf}. For
per-application design: each folder's \texttt{README.md} +
\texttt{INTEGRATION.md}, and
\texttt{applications/APPLICATIONS\_INTRO.pdf}.

\end{document}
